\begin{abstract}
Language models often struggle with mathematical reasoning, largely due to the inherent intricacies and formal structure characteristic of mathematical content. To address this, we present DeepSeekMath~7B, an open language model that extends the pre-training of DeepSeek-Coder-Base-v1.5 7B using an additional 120B tokens focused on mathematical data obtained from Common Crawl, as well as incorporating data from both natural language and code domains. DeepSeekMath~7B attains a notable 51.7\% on the highly challenging MATH benchmark at the competition level, all without utilizing external solvers or ensemble methods—narrowing the performance gap to leading models such as Gemini-Ultra and GPT-4. Moreover, generating 64 samples per problem for self-consistency yields a performance of 60.9\% on MATH. The strengths of DeepSeekMath~in mathematical reasoning result from two primary innovations: (1) the careful design of a data filtration and selection pipeline that effectively exploits a vast trove of web resources; and (2) the proposal of Group Relative Policy Optimization (GRPO), an adaptation of Proximal Policy Optimization (PPO), which enhances both mathematical reasoning and the efficiency of memory usage during reinforcement learning.
\end{abstract}

\section{Introduction}

Large language models (LLMs) have markedly transformed AI’s capacity for mathematical reasoning, driving progress on quantitative reasoning tasks \citep{MATH} as well as on benchmarks in geometry reasoning \citep{trinh2024solving}. Additionally, these models have demonstrated substantial value in providing support to humans tackling difficult mathematical challenges \citep{tao}. Nonetheless, leading-edge systems such as GPT-4 \citep{gpt4} and Gemini-Ultra \citep{gemini} remain inaccessible to the public, and current open-source models still significantly lag in terms of performance.

To address this gap, we present DeepSeekMath, a language model specialized for mathematics that raises the standard for open-source models and nearly rivals GPT-4 on established academic evaluations. This is accomplished by constructing the DeepSeekMath Corpus—a massive, high-quality pre-training dataset containing 120B mathematics-focused tokens. The corpus is assembled from Common Crawl (CC) through the use of a classifier based on fastText \citep{joulin2016fasttext}. The classifier is first trained using positive samples sourced from OpenWebMath \citep{openwebmath} and a variety of non-mathematical webpages as negatives. In the next stage, the trained classifier is deployed to mine more positive samples from the CC; these are then filtered and enhanced through human review. The classifier is retrained with this larger, cleaner dataset to optimize performance. Empirical results confirm the superior quality of the assembled corpus: our DeepSeekMath-Base 7B model attains 64.2\% accuracy on GSM8K \citep{gsm8k} and 36.2\% on the highly challenging MATH dataset \citep{MATH}, surpassing the performance achieved by Minerva 540B \citep{minerva}. Furthermore, thanks to its multilingual composition, the DeepSeekMath Corpus enables measurable gains on Chinese-language mathematical benchmarks \citep{wei2023cmath,agieval}. We see our work in mathematical data curation as an initial contribution for the research community, with abundant possibilities for future enhancements.

DeepSeekMath-Base is built upon DeepSeek-Coder-Base-v1.5 7B \citep{deepseek-coder}, given our finding that initializing from a model trained specifically for code yields superior results compared to leveraging a general-purpose LLM. In addition, we find that training on mathematical data not only strengthens the model's mathematical competencies, but also boosts its performance on the MMLU \citep{mmlu} and BBH \citep{bbh} benchmarks, thus supporting broader enhancements in reasoning skills.

Following the pre-training stage, we conduct mathematical instruction tuning on DeepSeekMath-Base utilizing datasets featuring chain-of-thought \citep{cot}, program-of-thought \citep{pot,pal}, and tool-integrated reasoning \citep{tora} strategies. The resultant DeepSeekMath-Instruct 7B model outperforms other 7B models and demonstrates performance on par with leading 70B instruction-tuned open-source models.

In addition, we propose Group Relative Policy Optimization (GRPO), a novel reinforcement learning (RL) approach derived from Proximal Policy Optimization (PPO) \citep{schulman2017proximal}. GRPO dispenses with the need for a critic network, instead utilizing group-based scores as its baseline, leading to substantial reductions in computational and training demands. Employing only a selected subset of English instruction tuning data, GRPO achieves notable gains compared to the already strong DeepSeekMath-Instruct model, yielding improvements across both in-domain (GSM8K: 82.9\% $\rightarrow$ 88.2\%, MATH: 46.8\% $\rightarrow$ 51.7\%) and out-of-domain tasks (for example, CMATH: 84.6\% $\rightarrow$ 88.8\%) during the RL phase. We further introduce a unified framework that encompasses and clarifies diverse approaches, such as Rejection Sampling Fine-Tuning (RFT) \citep{yuan2023scaling}, Direct Preference Optimization (DPO) \citep{dpo}, PPO, and our GRPO. This unified viewpoint reveals that these algorithms may all be interpreted as either direct or streamlined forms of RL. We carry out comprehensive empirical investigations—examining, for example, online versus offline training, outcome versus process-based supervision, and single-turn in contrast to iterative RL—enabling a deeper examination of the fundamental components underlying this unified approach. Finally, we provide an analysis detailing how our RL strategy enhances the performance of instruction-tuned models and outline promising avenues for developing more efficient RL algorithms within the framework of our unified paradigm.

\subsection{Contributions}
Our work contributes scalable mathematical pre-training and conducts a thorough investigation and evaluation of reinforcement learning techniques.

\textbf{Math Pre-Training at Scale}

\begin{itemize}[topsep=0pt]
    \item Our study demonstrates that publicly available Common Crawl data holds significant potential for mathematical applications. Through the deployment of a carefully crafted data selection framework, we assembled the DeepSeekMath Corpus: a curated dataset comprising 120B tokens sourced from web pages screened specifically for mathematical relevance. This corpus surpasses the dataset of math web pages utilized in Minerva \citep{minerva} by nearly sevenfold and is approximately nine times larger than the recently introduced OpenWebMath dataset \citep{openwebmath}.

    \item The DeepSeekMath-Base 7B model, once pre-trained, matches the performance of Minerva 540B \citep{minerva}, illustrating that scaling up parameter counts is not the sole determinant of mathematical reasoning prowess. Notably, smaller architectures trained on high-quality data can also exhibit strong mathematical problem-solving abilities.

    \item We present results from our experiments involving math-specific training regimens. Pre-training on code prior to undertaking math tasks substantially enhances the model's performance in mathematical problem-solving, effective both with and without auxiliary tools. These outcomes contribute insights into the debated topic of whether pre-training on code fosters improved reasoning, to which our evidence suggests an affirmative effect—at least within mathematical reasoning contexts.

    \item Contrary to widespread practice, the inclusion of arXiv papers in the training corpus—a common approach in various math-focused works—does not yield significant gains on any mathematical benchmarks considered in this research.
\end{itemize}

\textbf{Exploration and Analysis of Reinforcement Learning}

\begin{itemize}[topsep=0pt]
    \item We present Group Relative Policy Optimization (GRPO), a reinforcement learning algorithm characterized by its efficiency and efficacy. Unlike Proximal Policy Optimization (PPO), GRPO dispenses with the need for a critic model, instead leveraging group scores to estimate the baseline. This shift substantially lowers the computational resources required for training.
    \item Through our experiments, we show that GRPO substantially improves the performance of our instruction-tuned model, DeepSeekMath-Instruct, relying solely on instruction-tuning data. Notably, we also detect improvements in out-of-domain generalization during reinforcement learning with GRPO.
    \item We establish a comprehensive framework that brings coherence to various approaches including RFT, DPO, PPO, and GRPO. Our investigations encompass a series of rigorous empirical studies—comparing, for instance, online and offline training, outcome and process supervision, as well as single-turn versus iterative reinforcement learning—in order to scrutinize the core components of this framework.
    \item Drawing from insights afforded by our unified framework, we delve into the factors that drive the success of reinforcement learning. Furthermore, we distill several promising directions for advancing the reinforcement learning capabilities of LLMs.
\end{itemize}

\subsection{Summary of Evaluations and Metrics}

\begin{itemize}[topsep=0pt]
    \item \textbf{English and Chinese Mathematical Reasoning}: Our evaluation systematically examines model performance across both English and Chinese benchmark datasets, spanning mathematical questions from elementary through collegiate curricula. The English benchmarks used include GSM8K \citep{gsm8k}, MATH \citep{MATH}, SAT \citep{llemma}, OCW Courses \citep{minerva}, and MMLU-STEM \citep{mmlu}. For Chinese tasks, we utilize MGSM-zh \citep{mgsm}, CMATH \citep{wei2023cmath}, Gaokao-MathCloze \citep{agieval}, and Gaokao-MathQA \citep{agieval}. Our evaluation measures the models’ proficiency in generating fully self-contained written solutions in the absence of tools, as well as their capacity to solve questions using Python.

    DeepSeekMath-Base achieves competitive results against the proprietary Minerva 540B \citep{minerva} on English mathematical challenges, while substantially outperforming all openly available base models—including Mistral 7B \citep{mistral} and Llemma-34B \citep{llemma}—irrespective of whether math-specific pre-training has been conducted. Particularly in Chinese mathematics tasks, DeepSeekMath-Base delivers even stronger results, which we attribute to our deliberate inclusion of high-quality non-English mathematical data, contrasting with prior methodologies \citep{minerva,llemma} that use exclusively English data. Furthermore, once mathematical instruction tuning and reinforcement learning are applied, both DeepSeekMath-Instruct and DeepSeekMath-RL exhibit notable improvements, achieving, for the first time in the open-source community, over 50\% accuracy on the competitive MATH benchmark.

\item \textbf{Formal Mathematics}: We assess the capabilities of DeepSeekMath-Base on the informal-to-formal theorem proving task presented in \citep{dsp_proof}, utilizing miniF2F \citep{minif2f} and employing Isabelle \citep{isabelle} as the selected proof assistant. DeepSeekMath-Base achieves notable performance in few-shot autoformalization scenarios.

\item \textbf{Natural Language Understanding, Reasoning, and Code}: For a well-rounded evaluation of models' general comprehension, logical reasoning, and programming proficiency, we benchmark DeepSeekMath-Base using the Massive Multitask Language Understanding (MMLU) suite \citep{mmlu}, which includes 57 varied multiple-choice topics, the BIG-Bench Hard (BBH) collection \citep{bbh} featuring 23 tasks centered on complex, multi-step reasoning, and the HumanEval \citep{codex} and MBPP \citep{mbpp} tests, both commonly employed for assessing coding model performance. Pre-training on mathematical content is found to enhance both linguistic understanding and reasoning capabilities.

\section{Math Pre-Training}

\subsection{Data Collection and Decontamination} In this section, we detail the methodology for assembling the DeepSeekMath Corpus utilizing data sourced from Common Crawl. Illustrated in Figure \ref{fig:data_collect}, the pipeline presented here is iterative, illustrating a structured procedure for building a large-scale mathematical dataset from Common Crawl. The process is initiated with a seed dataset—comprising a relatively small but high-quality collection of mathematics-focused resources. This strategy is also transferable to other specialized areas, including programming.

To initiate the process, we select OpenWebMath \citep{openwebmath}, which comprises reputable mathematical content from the web, as our foundation seed corpus. Leveraging this dataset, we train a fastText classifier \citep{joulin2016fasttext} to retrieve additional web pages that resemble those in OpenWebMath. For training the model, we sample 500,000 items from the seed corpus as positive cases and 500,000 randomly chosen Common Crawl pages as negatives. The training is conducted using an open-source fastText implementation\footnote{\url{https://fasttext.cc}} with parameter settings as follows: vector size set to 256, learning rate at 0.1, a word n-gram maximum of 3, a minimum frequency threshold of 3 for words, and three training epochs. To handle redundancy, we apply URL-based deduplication as well as near-duplicate detection on Common Crawl, narrowing it down to a deduplicated set of 40B HTML pages. Utilizing the trained fastText model, we identify mathematical content within this deduplicated pool. To enhance the quality of selected documents, we assign ranks based on the fastText model’s output scores, retaining only the highest scoring items. The quantity of retained content is determined by conducting pre-training experiments on subsets containing the top 40B, 80B, 120B, and 160B tokens. In the initial round, we opt to preserve the segment comprising the top 40B tokens.

Following the initial round of data acquisition, a substantial number of mathematical web pages have yet to be gathered. This shortfall is largely attributable to the limited diversity within the positive examples on which the fastText model was originally trained. To address this, we augment the seed corpus by integrating more mathematical web sources, thereby allowing for enhanced model training. Our method begins by partitioning the full set of Common Crawl data into separate domains, each defined as the collection of web pages that share a common base URL. For each domain, we compute the proportion of web pages retrieved in the first iteration. Domains with over 10\% of their pages collected are marked as math-related (for instance, \url{mathoverflow.net}).

Next, we conduct manual annotation of URLs that are specifically associated with mathematical content within these math-related domains (such as \url{mathoverflow.net/questions}). Previously uncollected web pages that fall under these annotated URLs are appended to the seed set. By doing so, we accumulate a broader and richer collection of positive samples, which in turn helps us to further refine the fastText model and substantially improve recall of mathematical content in the next round of extraction. Upon completion of four total rounds, the process yields 35.5M mathematical web pages, collectively encompassing 120B tokens. During the fourth round, it becomes apparent that nearly 98\% of the identified pages had already been collected in the preceding round, leading us to conclude the data gathering effort at that stage.

In order to prevent benchmark contamination, our approach aligns with \cite{deepseek-coder} by eliminating web pages that include questions or answers derived from well-known English mathematical benchmarks such as GSM8K \citep{gsm8k} and MATH \citep{MATH}, as well as Chinese benchmarks like CMATH \citep{wei2023cmath} and AGIEval \citep{agieval}. The procedure involves discarding any portion of text in our mathematics training dataset if it contains a 10-gram sequence that perfectly matches any substring found in these evaluation benchmarks. For benchmark sources where the text length is below 10 grams but at least 3 grams, we utilize a precise string-matching method to identify and filter out web pages that may be contaminated.

\subsection{Validating the Quality of the DeepSeekMath~Corpus}

To assess the relative quality of the DeepSeekMath~Corpus, we conduct pre-training experiments that compare it to several notable recent math-focused corpora:

\begin{itemize}[topsep=0pt]
    \item \textbf{MathPile} \citep{mathpile}: This corpus contains 8.9B tokens and aggregates data from sources such as textbooks, Wikipedia, ProofWiki, CommonCrawl, StackExchange, and arXiv, with more than 85\% of its content originating from arXiv.
    \item \textbf{OpenWebMath} \citep{openwebmath}: Consisting of 13.6B tokens, OpenWebMath comprises mathematical data mined from CommonCrawl using targeted filtering methods.
    \item \textbf{Proof-Pile-2} \citep{llemma}: This dataset blends OpenWebMath, AlgebraicStack (which provides 10.3B tokens of math-related code), and arXiv documents (amounting to 28.0B tokens). For experiments with Proof-Pile-2, we adopt the arXiv:Web:Code ratio of 2:4:1, as specified in \cite{llemma}.
\end{itemize}

\subsubsection{Training Setting}
\label{sec:quality-policy}
We fine-tune a general-purpose pretrained language model containing 1.3B parameters—built on the same architecture as DeepSeek LLMs \citep{deepseek-llm}, referred to here as DeepSeek-LLM 1.3B—through math-specific training. For each mathematical dataset, a separate model instance is trained for a total of 150B tokens. The entirety of the training process is carried out using the lightweight and efficient HAI-LLM training framework \citep{haillm}. Adhering to the protocols established by DeepSeek LLMs, the AdamW optimizer \citep{adamW} is employed with hyperparameters $\beta_1=0.9$, $\beta_2=0.95$, and $\mathrm{weight\_decay}=0.1$. The learning rate is scheduled via a multi-step scheme: after an initial 2,000 warmup steps, it is increased to its peak value, then reduced to 31.6\% at 80\% completion of training, and further lowered to 10.0\% of the maximum for the last 10\% of training. The peak learning rate is fixed at 5.3e-4. Training is conducted with a batch size of 4M tokens, utilizing a context window of 4K tokens.

\subsubsection{Evaluation Results}

\textbf{The DeepSeekMath~Corpus distinguishes itself by its superior quality, extensive multilingual coverage, and unprecedented scale.}

\begin{itemize}[topsep=0pt]
    \item \textbf{High-quality}:
    We assess the corpus by analyzing model performance across eight downstream mathematical benchmarks utilizing few-shot chain-of-thought prompting \cite{cot}. Table \ref{tab:corpora_comparison} reveals that models trained with the DeepSeekMath~Corpus consistently achieve top results. Furthermore, Figure \ref{fig:corpora_comparisons} illustrates that, after exposure to 50B tokens (the equivalent of a single full epoch over Proof-Pile-2), the DeepSeekMath-trained model outperforms its Proof-Pile-2 counterpart, suggesting a higher average quality in the DeepSeekMath~Corpus.
    \item \textbf{Multilingual}:
    DeepSeekMath~Corpus comprises mathematical material spanning several languages, with English and Chinese constituting its primary linguistic bases. According to Table \ref{tab:corpora_comparison}, pretraining on the DeepSeekMath~Corpus yields improved mathematical reasoning abilities in both English and Chinese. By contrast, other available mathematical datasets focus predominantly on English, and do not contribute as much—and sometimes detract—from mathematical task performance in Chinese.
    \item \textbf{Large-scale}:
    The DeepSeekMath~Corpus greatly surpasses alternative mathematical corpora in size. As shown in Figure \ref{fig:corpora_comparisons}, DeepSeek-LLM 1.3B trained on DeepSeekMath~Corpus demonstrates a more rapid and sustained performance increase compared to models trained on much smaller baseline datasets—which have already been cycled through multiple times and consequently exhibit stagnation in further performance gains.
\end{itemize}

\subsection{Training and Evaluating DeepSeekMath-Base 7B}

This section details the development of DeepSeekMath-Base 7B, a foundational model characterized by advanced mathematical reasoning capabilities. The model is initialized using DeepSeek-Coder-Base-v1.5 7B \citep{deepseek-coder} as a starting point and is further trained on a dataset of 500B tokens. The composition of the training data consists of 56\% DeepSeekMath Corpus, 4\% AlgebraicStack, 10\% arXiv, 20\% Github code, and 10\% natural language texts sourced from the Common Crawl in both English and Chinese. The majority of training procedures align with those described in Section \ref{sec:quality-policy}, with the exceptions of utilizing a peak learning rate of 4.2e-4 and setting the batch size to 10M tokens.

We perform an in-depth evaluation of DeepSeekMath-Base 7B, examining its proficiency in independently generating full mathematical solutions without support from auxiliary tools, its competence in tackling mathematical tasks when tool assistance is available, and its aptitude for formal theorem verification. In addition to its mathematical abilities, we present a broader overview of the model's foundational capabilities, such as its competence in natural language comprehension, logical reasoning, and coding.

\paragraph{Mathematical Problem Solving with Step-by-Step Reasoning}

We assess DeepSeekMath-Base on its ability to address mathematical questions utilizing few-shot chain-of-thought prompting \citep{cot} across a set of eight standardized benchmarks in both English and Chinese. These assessments span quantitative reasoning datasets (such as GSM8K \citep{gsm8k}, MATH \citep{MATH}, and CMATH \citep{wei2023cmath}) as well as multiple-choice collections (like MMLU-STEM \citep{mmlu} and Gaokao-MathQA \citep{agieval}), encompassing a broad array of mathematical disciplines and covering topics from basic to collegiate-level complexity.

As detailed in Table \ref{tab:base_math_cot}, DeepSeekMath-Base 7B consistently achieves the highest scores on all eight evaluation benchmarks among open-source base models, including the extensively adopted Mistral 7B \citep{mistral} and the more recent Llemma 34B \citep{llemma}, which incorporates math-specific training from Proof-Pile-2 \citep{llemma}. Of particular note, on the challenging competition-grade MATH dataset, DeepSeekMath-Base exceeds the best-performing open-source alternatives by an absolute margin of more than 10\%, while also outperforming Minerva 540B \citep{minerva}—a proprietary model that, at 77 times the parameter count, leverages PaLM \citep{palm} and receives additional training on mathematical material.

\paragraph{Mathematical Problem Solving with Tool Use} For the assessment of program-assisted mathematical reasoning, we examine GSM8K and MATH with a few-shot program-of-thought prompt framework \citep{pot,pal}. Here, models tackle problems by generating Python scripts that may invoke libraries such as \textit{math} and \textit{sympy} to handle sophisticated calculations. The final solution is determined by the output of the executed script. Table \ref{tab:base_math_pot_proof} indicates that DeepSeekMath-Base 7B establishes a new performance benchmark, surpassing the previous state-of-the-art set by Llemma 34B.

\paragraph{Formal Mathematics}

The automation of formal proofs plays a critical role in verifying the correctness and robustness of mathematical arguments while improving workflow efficiency, and has garnered considerable interest in recent years. In this study, we assess the capabilities of DeepSeekMath-Base 7B on the informal-to-formal proof generation task introduced in \citep{dsp_proof}. This task involves constructing a formal proof using an informal statement, its formalized version, and an informal proof. Our experiments utilize the miniF2F benchmark \citep{minif2f}, which is designed to evaluate models on formalized Olympiad-level mathematics. We employ few-shot prompting to produce Isabelle-formalized proofs for each benchmark problem. Consistent with the methodology in \cite{dsp_proof}, our approach combines model-generated proof sketches with the automated Sledgehammer prover \citep{sledgehammer}, which completes the remaining proof details. According to the results displayed in Table \ref{tab:base_math_pot_proof}, DeepSeekMath-Base 7B achieves notable effectiveness in the autoformalization of mathematical proofs.

\paragraph{Natural Language Understanding, Reasoning, and Code}

To assess natural language understanding, we utilize the MMLU benchmark \citep{mmlu}; for reasoning abilities, we rely on BBH \citep{bbh}; and for coding proficiency, we employ HumanEval \citep{codex} alongside MBPP \citep{mbpp}. Table \ref{tab:understanding_reasoning_code} demonstrates that DeepSeekMath-Base 7B delivers substantial improvements over the previous DeepSeek-Coder-Base-v1.5 \citep{deepseek-coder} in both MMLU and BBH tasks, suggesting that training on mathematical data contributes positively to language understanding and reasoning. Furthermore, with the continued incorporation of code tokens during training, DeepSeekMath-Base 7B preserves the coding performance of DeepSeek-Coder-Base-v1.5 on HumanEval and MBPP. In aggregate, DeepSeekMath-Base 7B considerably surpasses the general-purpose Mistral 7B \citep{mistral} across the three benchmarks focused on reasoning and programming.

\section{Supervised Fine-Tuning}

\subsection{SFT Data Curation}

We develop a comprehensive mathematical instruction-tuning dataset comprising both English and Chinese problems drawn from a variety of mathematical disciplines and representing a range of difficulty levels. Each problem is accompanied by a corresponding solution articulated in one of several formats, namely chain-of-thought (CoT) \citep{cot}, program-of-thought (PoT) \citep{pot,pal}, or a tool-integrated reasoning approach \citep{tora}. In total, our dataset consists of 776K annotated training instances.

\begin{itemize}[topsep=0pt]
    \item \textbf{English mathematical datasets}: We enrich GSM8K and MATH datasets by providing tool-integrated solution annotations, incorporate a selection from MathInstruct \citep{MathInstruct}, and use the training subset of Lila-OOD \citep{lila} which features problems resolved using either CoT or PoT techniques. The English portion encompasses a broad array of mathematical subjects, including but not limited to algebra, probability, number theory, calculus, and geometry.
    \item \textbf{Chinese mathematical datasets}: We assemble a collection of Chinese K-12 mathematics problems distributed across 76 subdomains, for example, linear equations. Solutions to these problems are annotated in both CoT and tool-integrated reasoning formats.
\end{itemize}

\subsection{Training and Evaluating DeepSeekMath-Instruct 7B}

This section presents DeepSeekMath-Instruct 7B, developed through mathematical instruction fine-tuning of DeepSeekMath-Base. During training, samples are concatenated at random until the 4K-token maximum context limit is reached. The fine-tuning process comprises 500 steps, utilizing a batch size of 256 and a fixed learning rate set to 5e-5.

The mathematical capabilities of the models are assessed under two settings—with and without tool assistance—across four quantitative reasoning benchmarks in both English and Chinese. We compare DeepSeekMath-Instruct 7B to state-of-the-art models current at the time of writing:

\begin{itemize}[topsep=0pt]
    \item \textbf{Closed-source models} considered for comparison include: (1) the GPT series, of which GPT-4 \citep{gpt4} and GPT-4 Code Interpreter~\footnote{\url{https://openai.com/blog/chatgpt-plugins\#\#code-interpreter}} represent the top-performing options, (2) Gemini Ultra and Pro \citep{gemini}, (3) Inflection-2 \citep{inflection-2}, (4) Grok-1~\footnote{\url{https://x.ai/model-card}}, and recent releases from Chinese tech companies, including (5) Baichuan-3~\footnote{\url{https://www.baichuan-ai.com}}, and (6) the most recent GLM-4~\footnote{\url{https://open.bigmodel.cn/dev/api\#glm-4}}

    from the GLM family \citep{glm}.

These models are primarily intended for broad applications, and the majority have completed multiple stages of alignment.

\item \textbf{Open-source models} consist of both general-purpose and math-optimized systems, including: (1) DeepSeek-LLM-Chat 67B \citep{deepseek-llm}, (2) Qwen 72B \citep{qwen}, (3) SeaLLM-v2 7B \citep{seallm}, and (4) ChatGLM3 6B \citep{chatglm3} among the general models; and a suite of math-focused enhancements such as (5) InternLM2-Math 20B\footnote{\url{https://github.com/InternLM/InternLM-Math}}, derived from InternLM2 with dedicated math training and subsequent instruction tuning, (6) Math-Shepherd-Mistral 7B, which introduces PPO training \citep{schulman2017proximal} with a process-supervised reward framework applied to Mistral 7B \citep{mistral}, (7) the WizardMath suite \citep{wizardmath}, which augments mathematical reasoning capabilities in both Mistral 7B and Llama-2 70B \citep{llama2} via evolve-instruct (an AI-evolved instruction tuning method) alongside PPO training and utilizing datasets largely sourced from GSM8K and MATH, (8) MetaMath 70B \citep{metamath}, wherein Llama-2 70B is fine-tuned using expanded versions of GSM8K and MATH, (9) ToRA 34B \cite{tora}, created by fine-tuning CodeLlama 34B to facilitate mathematical reasoning integrated with external tools, and (10) MAmmoTH 70B \citep{MathInstruct}, realized by instruction-tuning Llama-2 70B on MathInstruct.

According to the results presented in Table \ref{tab:sft_rl_math}, when evaluation restricts access to external tools, DeepSeekMath-Instruct 7B exhibits strong capabilities in conducting step-by-step reasoning. On the competition-level MATH dataset, our model outperforms all other open-source models, as well as most closed-source systems (such as Inflection-2 and Gemini Pro), achieving at least a 9\% absolute improvement. This performance holds even in comparison to much larger models (e.g., Qwen 72B) or those that have been fine-tuned using math-specific reinforcement learning techniques (such as WizardMath-v1.1 7B). Although DeepSeekMath-Instruct achieves results comparable to Chinese proprietary models like GLM-4 and Baichuan-3 on the MATH benchmark, its performance still falls short of leading models like GPT-4 and Gemini Ultra.

When assessed in a context that permits the utilization of both natural language reasoning and programmatic tools for solving problems, DeepSeekMath-Instruct 7B nearly achieves a 60\% accuracy rate on the MATH benchmark, outperforming all other open-source competitors. For the remaining evaluation sets, our model demonstrates performance on par with DeepSeek-LLM-Chat 67B, the former leading model that is an order of magnitude larger in size.

\section{Reinforcement Learning}

\subsection{Group Relative Policy Optimization} Reinforcement learning (RL) has demonstrated its efficacy in augmenting the mathematical reasoning skills of large language models (LLMs) beyond what is achieved in the Supervised Fine-Tuning (SFT) phase \citep{wang2023math,wizardmath}. Here, we present Group Relative Policy Optimization (GRPO), a reinforcement learning technique that is both efficient and robust.

\subsubsection{From PPO to GRPO} Proximal Policy Optimization (PPO) \citep{schulman2017proximal} remains one of the most prominent actor-critic RL algorithms and is extensively adopted in the RL fine-tuning process for LLMs \citep{ouyang2022training}. PPO enhances LLMs by seeking to maximize a particular surrogate objective, formalized as:

\begin{equation}
\footnotesize
    \mathcal{J}_{PPO}(\theta) = \mathbb{E}{[q \sim P(Q), o \sim \pi_{\theta_{old}}(O|q)]} \frac{1}{|o|} \sum_{t=1}^{|o|} \min \left[ \frac{\pi_\theta(o_{t} | q, o_{<t})}{\pi_{\theta_{old}}(o_{t} | q, o_{<t})} A_{t}, \text{clip} \left( \frac{\pi_\theta(o_{t} | q, o_{<t})}{\pi_{\theta_{old}}(o_{t} | q, o_{<t})}, 1 - \varepsilon, 1 + \varepsilon \right)  A_{t} \right] ,
\end{equation}

Here, $\pi_{\theta}$ denotes the updated policy model, while $\pi_{\theta_{old}}$ refers to the policy prior to the latest update. The variables $q$ and $o$ correspond to questions and outputs, which are respectively drawn from the question dataset and generated by the previous policy $\pi_{\theta_{old}}$. The hyper-parameter $\varepsilon$, which is related to clipping, is employed in PPO to maintain stable training dynamics. The advantage term, $A_t$, is determined using the Generalized Advantage Estimation (GAE) method \citep{gae}, leveraging both the reward sequence $\{r_{\ge t}\}$ and a value function $V_{\psi}$ that is trained in parallel with the policy model. To prevent excessive exploitation of the reward model, it is standard to include a per-token Kullback-Leibler (KL) divergence penalty between the policy and a reference model in the reward computation at each step \citep{ouyang2022training}, given by

\begin{equation}
    r_{t} = r_\varphi(q, o_{\le t}) - \beta \log\frac{\pi_{\theta}(o_{t}|q, o_{<t})}{\pi_{ref}(o_{t}|q, o_{<t})},
\label{eq:PPO-reward}
\end{equation}

where the reward model is denoted by $r_\varphi$, the reference model $\pi_{ref}$ is typically the original SFT model, and $\beta$ indicates the magnitude of the KL regularization term.

In PPO, the value function is often represented by a model that is roughly as large as the policy network itself, resulting in notable computational and memory demands. Furthermore, throughout the reinforcement learning process, the value function acts as a baseline to compute the advantage, thereby aiming to reduce variance. In large language model (LLM) applications, however, it is common for the reward model to assign a reward only to the final token of a generated sequence, which poses challenges when attempting to train a value function capable of providing accurate value estimations at each token position. To circumvent these challenges, and as illustrated in Figure \ref{fig:grpo}, we introduce Group Relative Policy Optimization (GRPO). GRPO eliminates the requirement for an explicit value function—unlike PPO—by leveraging the mean reward of a set of responses, all generated for the same prompt, as the baseline. Concretely, for any question $q$, GRPO generates a set of outputs $\{o_1, o_2, \cdots, o_G\}$ using the previous policy $\pi_{\theta_{old}}$, and then updates the policy model by maximizing the subsequent objective:

\begin{equation}
\footnotesize
\begin{split}
    \mathcal{J}_{GRPO}(\theta) &= \mathbb{E}{[q \sim P(Q), \{o_i\}_{i=1}^G \sim \pi_{\theta_{old}}(O|q)]}  \\
    & \frac{1}{G}\sum_{i=1}^G\frac{1}{|o_i|} \sum_{t=1}^{|o_i|} \left\{ \min \left[ \frac{\pi_\theta(o_{i,t} | q, o_{i,<t})}{\pi_{\theta_{old}}(o_{i,t} | q, o_{i,<t})} \hat{A}_{i,t}, \text{clip} \left( \frac{\pi_\theta(o_{i,t} | q, o_{i,<t})}{\pi_{\theta_{old}}(o_{i,t} | q, o_{i,<t})}, 1 - \varepsilon, 1 + \varepsilon \right)  \hat{A}_{i,t} \right] - \beta \mathbb{D}_{KL}\left[\pi_{\theta} || \pi_{ref}\right]\right\} ,
\end{split}
\label{eq:GRPO-obj}
\end{equation}

Here, $\varepsilon$ and $\beta$ denote hyper-parameters, and $\hat{A}_{i,t}$ stands for the advantage, which is computed by considering the relative rewards among the outputs within each group; a more thorough explanation will be provided in the following subsections. By constructing advantages in this group-relative manner, GRPO capitalizes on the inherently comparative design of reward models, which are generally developed on pairwise output evaluations for the same prompt. Furthermore, rather than integrating the KL penalty directly into the reward, GRPO enforces regularization by explicitly including the KL divergence term between the learned and the reference policy within the loss function. This formulation simplifies the computation of $\hat{A}_{i,t}$ by avoiding its entanglement with the KL penalty. In contrast to the KL penalty approach in (\ref{eq:PPO-reward}), GRPO adopts the following unbiased estimator for the KL divergence \citep{kl_approx}:

\begin{equation}
\small
    \mathbb{D}_{KL}\left[\pi_{\theta} || \pi_{ref}\right] = \frac{\pi_{ref}(o_{i,t}|q,o_{i,<t})}{\pi_{\theta}(o_{i,t}|q,o_{i,<t})}- \log\frac{\pi_{ref}(o_{i,t}|q,o_{i,<t})}{\pi_{\theta}(o_{i,t}|q,o_{i,<t})} - 1,
\end{equation}

which is guaranteed to be positive.

\begin{algorithm}[t]
  \small
  \caption{Iterative Group Relative Policy Optimization}
  \textbf{Input} Starting policy model $\pi_{\theta_{\text{init}}}$; reward models $r_\varphi$; task dataset $\mathcal{D}$; 
  hyperparameters $\varepsilon$, $\beta$, $\mu$
  \begin{algorithmic}[1]
    \State Initialize policy model: $\pi_\theta \leftarrow \pi_{\theta_{\text{init}}}$
    \For{iteration = 1, \dots, I}
       \State Set reference model $\pi_{ref} \leftarrow \pi_{\theta}$
      \For{step = 1, \dots, M}
      \State Draw a batch $\mathcal{D}_b$ from $\mathcal{D}$
      \State Save the current policy model $\pi_{\theta_{old}} \leftarrow \pi_{\theta}$ 
      \State For each query $q$ in $\mathcal{D}_b$, generate $G$ samples $\{o_i\}_{i=1}^G \sim \pi_{\theta_{old}} (\cdot \mid q)$
      \State Evaluate rewards $\{r_i\}_{i=1}^{G}$ for all generated outputs $o_i$ using $r_{\varphi}$
      \State Estimate group-relative advantages $\hat{A}_{i,t}$ for each token position $t$ of each output $o_i$
      \For{GRPO iteration = 1, \dots, $\mu$}
        \State Optimize the policy model $\pi_{\theta}$ by increasing the GRPO objective (see Equation \ref{eq:GC-GRPO})
      \EndFor
    \EndFor
    \State Perform continual training of $r_\varphi$ using a replay buffer mechanism.
    \EndFor
  \end{algorithmic}
  \textbf{Output} $\pi_\theta$
  \label{alg:iter-grpo}
\end{algorithm}

\subsubsection{Outcome Supervision RL with GRPO} Specifically, given a question $q$, a set of output candidates $\{o_1, o_2, \cdots, o_G\}$ is generated by sampling from the previous policy, $\pi_{\theta_{old}}$. Each of these outputs is evaluated by a reward model, resulting in a set of rewards $\mathbf{r}=\{r_1, r_2, \cdots, r_G\}$. The obtained rewards are standardized by removing the group mean and scaling by the group's standard deviation. Within the outcome supervision paradigm, each sequence $o_i$ receives its final, normalized reward as the per-token advantage; thus, for all tokens in $o_i$, the advantage is $\hat{A}_{i, t} = \widetilde{r}_i = \frac{r_i- {\rm mean}(\mathbf{r})}{{\rm std}(\mathbf{r})}$. The policy is then improved by maximizing the objective specified in equation (\ref{eq:GRPO-obj}).

\subsubsection{Process Supervision RL with GRPO} Traditional outcome supervision delivers feedback solely at the conclusion of each output, which can be inadequate and inefficient for guiding the policy, particularly on complex mathematical problems. In line with \cite{wang2023math}, we investigate process supervision, a scheme that assigns rewards after every intermediate reasoning step. Concretely, for a given question $q$ and a collection of $G$ sampled outputs $\{o_1, o_2, \cdots, o_G\}$, a process reward model evaluates each reasoning step, assigning a sequence of rewards as follows: $\mathbf{R} = \{ \{r_1^{index(1)},\cdots,r_1^{index(K_1)}\}, \cdots,  \{r_G^{index(1)},\cdots,r_G^{index(K_G)}\} \}$. Here, $index(j)$ denotes the token index at which the $j$-th step concludes, and $K_i$ is the step count in the $i$-th output. To ensure consistency across rewards, we apply normalization based on their empirical mean and standard deviation, as: $\widetilde{r}_i^{index(j)} = \frac{r_i^{index(j)} - {\rm mean(\mathbf{R})}}{{\rm std(\mathbf{R})}}$.

Process supervision then computes the advantage for each token position by accumulating the normalized rewards of all subsequent reasoning steps, that is, $\hat{A}_{i, t} = \sum_{index(j) \ge t} \widetilde{r}_i^{index(j)}$. This advantage is subsequently utilized to optimize the policy with respect to the objective function specified in equation (\ref{eq:GRPO-obj}).

\subsubsection{Iterative RL with GRPO} During the reinforcement learning phase, the previously trained reward model may become inadequate to effectively guide the evolving policy model. To address this, we investigate iterative RL using GRPO. As depicted in Algorithm \ref{alg:iter-grpo}, the iterative GRPO framework involves generating fresh datasets for reward model training utilizing newly sampled outputs from the current policy model. The reward model is updated in an ongoing manner through a replay strategy that retains 10\% of prior training data. Following this, the policy model is established as the new reference model, and further training of the policy model proceeds under the guidance of the updated reward model.

\subsection{Training and Evaluating DeepSeekMath-RL}

Our RL experiments are performed utilizing DeepSeekMath-Instruct 7B. The RL training corpus consists of chain-of-thought questions associated with GSM8K and MATH, drawn from the SFT data and totaling approximately 144K examples. To focus on evaluating RL's effects on tasks underrepresented during the RL process, we exclude remaining SFT-derived questions. The procedure to assemble the reward model’s training dataset is adopted from \citep{wang2023math}. The reward model is initially fine-tuned on DeepSeekMath-Base 7B using a learning rate of 2e-5. In the GRPO setting, the learning rate for the policy model is set at 1e-6, with a KL penalty coefficient of 0.04. For each individual question, we generate $64$ candidate completions. The sequence length is capped at 1024 tokens, and each training batch comprises 1024 samples. The policy network receives a single update per exploration cycle. DeepSeekMath-RL 7B is evaluated on the same suite of benchmarks as DeepSeekMath-Instruct 7B. In this context, tasks like GSM8K and MATH—both requiring chain-of-thought reasoning—are classified as in-domain, while all remaining benchmarks are treated as out-of-domain.

Table \ref{tab:sft_rl_math} presents a comparison of the capabilities of open- and closed-source models employing both chain-of-thought and tool-augmented reasoning across English and Chinese evaluation sets. The results highlight several key observations: 1) DeepSeekMath-RL 7B achieves accuracy scores of 88.2\% on GSM8K and 51.7\% on MATH with chain-of-thought reasoning, outperforming every open-source model within the 7B to 70B parameter range, and even exceeding the performance of most closed-source models. 2) Notably, DeepSeekMath-RL 7B's training is limited exclusively to instruction tuning datasets in chain-of-thought format from GSM8K and MATH, with its initialization derived from DeepSeekMath-Instruct 7B. Although it is exposed to a narrower range of training examples, DeepSeekMath-RL 7B surpasses DeepSeekMath-Instruct 7B in all assessment metrics, underscoring the advantages conferred by reinforcement learning.

\section{Discussion}
This section provides an overview of our insights gained from pre-training and reinforcement learning experiments.

\subsection{Lessons Learnt in Pre-Training}

We begin by describing our experiences during pre-training. Unless stated otherwise, the training configurations specified in Section \ref{sec:quality-policy} will be used throughout. Additionally, it is important to mention that the DeepSeekMath Corpus discussed here utilizes an 89B-token collection produced during the second round of dataset assembly.

\subsubsection{Code Training Benefits Mathematical Reasoning}

A widely held, albeit largely untested, hypothesis posits that training with code enhances reasoning capabilities. Here, we present preliminary evidence specific to the mathematical domain: exposure to code during training strengthens models’ mathematical reasoning skills, regardless of whether external tools are employed.

To investigate the impact of code training on mathematical reasoning, we conducted experiments employing both a two-stage training procedure and a single-stage training baseline:

\textbf{Two-Stage Training}

\begin{itemize}[topsep=0pt]
    \item \textbf{Code Training for 400B Tokens $\rightarrow$ Math Training for 150B Tokens}: The DeepSeek-LLM 1.3B model is first trained on a dataset comprising 400B code tokens, after which it is further trained on an additional 150B math tokens;
    \item \textbf{General Training for 400B Tokens $\rightarrow$ Math Training for 150B Tokens}: For comparison, we conduct a control experiment in which the initial 400B tokens are drawn from a broad general-purpose corpus curated by DeepSeek-AI, replacing code tokens with general tokens, to assess whether pretraining on code specifically yields greater improvements in mathematical reasoning relative to a general domain pretraining.
\end{itemize}

\textbf{One-Stage Training}

\begin{itemize}[topsep=0pt]
    \item \textbf{Math Training on 150B Tokens}: The DeepSeek-LLM 1.3B model undergoes training with 150B math-specific tokens.
    \item \textbf{Combined Training with 400B Code Tokens and 150B Math Tokens}: Sequential math training after code pretraining has been found to decrease code capability. To address this, we analyze whether jointly training with a mixture of code and math tokens in a single phase not only boosts mathematical reasoning skills but also reduces catastrophic forgetting.
\end{itemize}

\paragraph{Results}
The outcomes detailed in Table \ref{tab:code-math} and Table \ref{tab:code-math-general-eval} compare downstream performance across various training regimes.

Training on code tokens contributes significantly to the improvement of program-based mathematical reasoning under both single-phase and two-phase approaches. Specifically, as reported in Table \ref{tab:code-math}, two-stage training with code-only tokens initially results in substantial gains on problems like GSM8K and MATH, solvable with Python tools. Subsequent training on math data in a second phase continues to enhance these results. Notably, one-stage training, involving a mixture of code and math tokens, both prevents the catastrophic forgetting observed in the sequential approach and leads to joint improvements for coding tasks (see Table \ref{tab:code-math-general-eval}) as well as program-aided mathematical reasoning (see Table \ref{tab:code-math}).

Training with code also enhances mathematical reasoning abilities in the absence of tool assistance. Within a two-phase training framework, the preliminary phase that focuses on code delivers notable yet moderate improvements. Furthermore, this stage facilitates more effective mathematical training in the subsequent phase, ultimately yielding superior performance. In contrast, merging code and mathematical tokens for concurrent, single-stage training leads to diminished mathematical reasoning when tool use is not permitted. A plausible explanation is that DeepSeek-LLM 1.3B, owing to its comparatively small size, is unable to accommodate both code and mathematical inputs concurrently to a satisfactory extent.

\subsubsection{ArXiv Papers Appear Unproductive for Mathematical Reasoning Gains}

Although arXiv papers are routinely incorporated in mathematical pre-training datasets \citep{minerva,gpt-f,llemma,mathpile}, systematic evaluation of their role in enhancing mathematical reasoning is scarce. Interestingly, our results indicate that arXiv papers do not meaningfully support the advancement of mathematical reasoning, contrary to what one might expect. Our investigation includes models of multiple sizes such as DeepSeek-LLM 1.3B and DeepSeek-Coder-Base-v1.5 7B \citep{deepseek-coder}, and utilizes arXiv-based datasets processed via distinct strategies:

\begin{itemize}[topsep=0pt]
    \item \textbf{MathPile} \citep{mathpile}: a dataset encompassing 8.9B tokens, curated by applying a range of cleaning and filtering heuristics, in which scientific arXiv documents constitute more than 85\% of the material;
    \item \textbf{ArXiv-RedPajama} \citep{redpajama}: this dataset comprises all arXiv LaTeX source files stripped of preambles, comments, macro definitions, and bibliographic entries, resulting in a collection with 28.0B tokens.
\end{itemize}

During our experimental evaluation, we conducted isolated training sessions, using DeepSeek-LLM 1.3B with 150B tokens and DeepSeek-Coder-Base-v1.5 7B with 40B tokens for each arXiv-specific dataset. The results suggest that incorporating arXiv papers does not substantially enhance mathematical reasoning abilities. When these models were trained solely on the arXiv dataset, there were either negligible gains or, in some cases, declines observed across an array of mathematical evaluation tasks with varying difficulty, as assessed in this work. These tasks encompassed quantitative reasoning benchmarks such as GSM8K and MATH (see Table \ref{tab:arxiv-cot}), multiple-choice tasks exemplified by MMLU-STEM (see Table \ref{tab:arxiv-cot}), as well as datasets grounded in formal mathematics such as miniF2F (see Table \ref{tab:arxiv_atp}).

Nonetheless, this finding warrants a cautious interpretation due to several limitations in our study scope. Specifically, we have not yet analyzed:

\begin{itemize}[topsep=0pt]
    \item The influence of arXiv-sourced data on distinct mathematical activities that are outside the current assessment, such as informalization—the translation of formal theorems or proofs into informal descriptions;
    \item The outcomes arising from mixing arXiv data with other forms of training data;
    \item Whether models of greater scale might be better able to leverage potential advantages of arXiv-derived text.
\end{itemize}

Accordingly, additional investigations will be essential, and we defer these for future research initiatives.

\subsection{Insights of Reinforcement Learning}
\subsubsection{Towards to a Unified Paradigm} Here, we aim to establish a generalized analytical framework for contrasting several training algorithms, including SFT, RFT, DPO, PPO, and GRPO, and proceed to experimentally investigate critical aspects of this unified paradigm. In its general form, the gradient of a training objective with respect to parameters $\theta$ can be expressed as follows:

\begin{equation}
    \nabla_{\theta}\mathcal{J}_{\textcolor{red}{\mathcal{A}}}(\theta) = \mathbb{E}[\underbrace{(q,o) \sim \textcolor{red}{\mathcal{D}}}_{Data \ Source}]\left( \frac{1}{|o|} \sum_{t=1}^{|o|}  \underbrace{GC_{{\mathcal{A}}}(q, o, t, \textcolor{red}{\pi_{{rf}}})}_{Gradient \ Coefficient}  \nabla_{\theta}\log \pi_{\theta}(o_t | q, o_{<t})\right).
\label{eq:objective}
\end{equation}

This formulation involves three principal elements: 1) the \textit{Data Source} $\mathcal{D}$, which specifies the dataset employed for training; 2) the \textit{Reward Function} $\pi_{{rf}}$, providing the rewards utilized as feedback during training; and 3) the \textit{Algorithm} $\mathcal{A}$, which integrates the collected training instances and corresponding reward feedback into the gradient coefficient $GC$, modulating the extent of penalization or reinforcement applied to the data. Building on this general framework, we examine several illustrative methodologies:

\begin{itemize}
    \item \textbf{Supervised Fine-tuning (SFT)}: In SFT, a pretrained model is further trained on a dataset curated by humans specifically for this phase.
    \item \textbf{Rejection Sampling Fine-tuning (RFT)}: RFT continues the training of the SFT model by utilizing a subset of its responses to SFT prompts, selected by filtering based on answer accuracy.
    \item \textbf{Direct Preference Optimization (DPO)}: DPO advances the SFT model by employing pairwise DPO loss to fine-tune it on a set of outputs that have been augmented and sampled from the SFT model.
    \item \textbf{Online Rejection Sampling Fine-tuning (Online RFT)}: Unlike traditional RFT, Online RFT starts with the SFT-initialized policy model and iteratively improves it by fine-tuning on freshly generated samples from the current policy itself.
    \item \textbf{PPO/GRPO}: These methods begin by initializing the policy with the SFT model, and then strengthen it using examples sampled from the evolving, real-time policy.
\end{itemize}

A summary of these methods' components is presented in Table \ref{tab:data_policy}. For a comprehensive explanation of the derivation process, please consult Appendix \ref{app:analysis-rl}.

\paragraph{Observation about Data Source} We categorize data sources into online sampling and offline sampling. In online sampling, training data is generated from real-time explorations conducted by the current training policy model. Conversely, offline sampling refers to using data sampled from the outputs of the initial SFT model. Methods such as RFT and DPO employ the offline sampling approach, whereas Online RFT and GRPO utilize online sampling.

As illustrated in Figure \ref{fig:rl-analysis}, our results indicate that Online RFT achieves substantially better performance than RFT across both benchmarks. During the early phases of training, Online RFT performs similarly to RFT; however, as training progresses, Online RFT exhibits clear superiority, underscoring the benefits of adopting an online training paradigm. This outcome aligns with expectations: initially, the actor and the SFT model remain quite similar, resulting in the collected data reflecting only slight variations. Over time, as training advances, the data generated by the actor diverges more markedly, making real-time data acquisition increasingly advantageous.

\paragraph{Observation about Gradient Coefficient} The algorithm utilizes the reward signal to determine the gradient coefficient, which in turn informs the update of model parameters. In our experimental framework, we distinguish between two forms of reward functions: `Rule' and `Model'. The `Rule' approach evaluates the appropriateness of a response by assessing answer correctness, whereas the `Model' strategy involves training a reward model to assign scores to each response, using data annotated through rule-based assessments. As indicated in Equations \ref{eq:GC-RFT} and \ref{eq:GC-GRPO}, there exists an important distinction between GRPO and Online RFT: GRPO exclusively modifies its gradient coefficient as a function of the reward value assigned by the reward model. This mechanism enables GRPO to differentially strengthen or weaken responses according to the degree of reward they receive. On the other hand, Online RFT does not share this adaptive property; it refrains from penalizing wrong answers and consistently applies the same degree of positive reinforcement to all correctly answered responses.

As illustrated in Figure \ref{fig:rl-analysis}, GRPO outperforms the online RFT approach, underscoring the effectiveness of modulating the coefficients for positive and negative gradients. Moreover, GRPO+PS delivers better results than GRPO+OS, demonstrating the advantages provided by the use of fine-grained, step-specific gradient adjustments. We also assess the impact of iterative RL, performing two consecutive rounds in our experiments. Figure \ref{fig:iter-rl} reveals that iterative RL yields notable performance gains, with the most pronounced improvements occurring after the initial iteration.

\subsubsection{Why RL Works?} In this study, we apply reinforcement learning techniques to a selected subset of instruction tuning data, achieving a marked increase in performance over the base instruction tuning model. To further investigate the mechanism underlying this improvement, we analyze the Pass@K and Maj@K accuracy metrics for both Instruct and RL-trained models across two benchmark datasets. According to Figure \ref{fig:maj-pass}, RL particularly enhances Maj@K accuracy, while Pass@K remains unaffected. This suggests that reinforcement learning chiefly bolsters the robustness of the model’s output distribution, \textbf{improving the likelihood that the correct answer is chosen from among the TopK responses, rather than fundamentally advancing model capability.} Likewise, \citep{wang2023making} observed a \textbf{misalignment problem} affecting reasoning tasks in SFT models, and further demonstrated that the reasoning skills of SFT models could be augmented by employing various preference alignment techniques \citep{yuan2023rrhf,song2023preference,wang2023making}.

\subsubsection{How to Achieve More Effective RL?}
\label{sec:effec_rl}
Our results indicate that RL demonstrates considerable effectiveness when applied to mathematical reasoning challenges. We introduce a unified framework that provides insight into a range of prominent training strategies. Under this scheme, these approaches can all be regarded as variations of either direct or abstracted RL methods. As encapsulated by Equation \ref{eq:objective}, there are three central elements to consider: Data Source, Algorithm, and Reward Function. Below, we discuss possible avenues for advancing each of these components.

\paragraph{Data Source} The data source serves as the foundational input for every training approach. When discussing RL, we specifically define the data source as the set of unlabeled questions alongside the responses produced by sampling from the policy model. In our experiments, we rely solely on questions taken from the instruction tuning phase and adopt a straightforward nucleus sampling method for output generation. We attribute the limited gains observed in our RL pipeline, primarily in Maj@K metrics, to this design choice. Future investigations will seek to deploy our RL framework on question prompts beyond the training distribution and to integrate \textbf{enhanced sampling (decoding) techniques}—such as those employing tree-search strategies \citep{yao2023tree}. Additionally, recent advancements in \textbf{efficient inference methods} \citep{xia-etal-2023-speculative,leviathan2023fast,kwon2023efficient,xia2024unlocking}, which critically influence the policy model's exploration capabilities, warrant substantial attention.

\paragraph{Algorithms} In the process of model parameter optimization, algorithms utilize the reward feedback to calculate gradient coefficients for updates. As indicated in Equation \ref{eq:objective}, current approaches generally place complete \textbf{TRUST} in the feedback generated by the reward function, leveraging it to adjust the conditional probability of tokens upward or downward. Nevertheless, this reward signal cannot be assumed to be consistently trustworthy, especially in tasks that are highly intricate. As an illustration, the PRM800K datasets \citep{lightman2023let}, despite the rigorous curation and expertise of annotators, still exhibit a misannotation rate of about 20\%\footnote{\url{https://github.com/openai/prm800k/issues/12\#issuecomment-1728491852}}. Thus, there is a need to investigate reinforcement learning algorithms that can effectively cope with noisy or imperfect reward signals. We contend that adopting \textbf{WEAK-TO-STRONG} alignment techniques \citep{burns2023weak} could fundamentally transform learning algorithm design.

\paragraph{Reward Function} The reward function serves as the primary driver of learning in RL settings, where it is typically represented by a neural reward model. We identify three key areas of focus for advancing reward models: 1) \textbf{Enhancing the generalization capabilities of the reward model.} For reinforcement learning to meaningfully improve large language models (LLMs), reward models must generalize effectively to out-of-distribution queries and sophisticated generated outputs; failure in this regard could lead to reinforcement learning merely consolidating pre-existing distributions without enhancing underlying competencies. 2) \textbf{Capturing and representing reward model uncertainty.} Expressing uncertainty can serve as a crucial connector between less reliable reward functions and learning algorithms that transition from weaker to more robust performance. 3) \textbf{Efficient development of high-fidelity process reward models} capable of delivering detailed, stepwise feedback to better support reasoning tasks \citep{lightman2023let,wang2023math}.

\section{Conclusion, Limitation, and Future Work} This work introduces DeepSeekMath, an open-source model that surpasses existing counterparts on the competition-grade MATH benchmark, narrowing the gap with leading proprietary systems. DeepSeekMath~builds upon DeepSeek-Coder-v1.5 7B, undergoing extended continued pre-training for 500B tokens, with 120B mathematical tokens from Common Crawl constituting a substantial portion of the dataset. Through comprehensive ablation analysis, our findings indicate that web pages are a rich resource for high-quality mathematical content, while arXiv's contributions may be less impactful than anticipated. We present Group Relative Policy Optimization (GRPO), an adaptation of Proximal Policy Optimization (PPO), which significantly enhances mathematical reasoning capabilities while reducing memory demands. Empirical evidence demonstrates that GRPO delivers further improvements even when DeepSeekMath-Instruct 7B attains top-tier benchmark scores. Additionally, we offer a unified interpretative framework for related methodologies and delineate promising avenues for advancing reinforcement learning strategies.

While DeepSeekMath~demonstrates strong results on benchmarks centered around quantitative reasoning, it tends to underperform relative to proprietary systems when it comes to geometry and theorem proving. For example, our initial evaluations reveal that the model is unable to address certain questions involving triangles and ellipses, suggesting the possibility of selection bias in the datasets used during pre-training and fine-tuning. Furthermore, the model's size limits its few-shot learning capabilities; whereas GPT-4 exhibits marked improvement when given few-shot examples, DeepSeekMath~shows little difference between zero-shot and few-shot settings. Moving forward, we aim to refine our curated data selection procedures to assemble a more robust and diverse pre-training corpus. Additionally, we will investigate avenues for enhancing reinforcement learning methods in LLMs as outlined in Section \ref{sec:effec_rl}.

\end{document}